%============================================================ 3. 문제 정의
\section{문제 정의와 평가 지표}
\label{sec:problem}

\subsection{교전 시나리오}
\label{sec:scenario}
% Table 1·2 는 §3 첫머리(p3)에서 선언한다. 선언 순서 scenario → metrics 가 인용 순서와
% 같아야 번호가 맞는다. Fig.2 는 p4 상단, Table 2 는 p5 상단에 놓인다(2026-09-16 실측).
% Table 2 를 p4 하단에 넣는 시도([!b], 선언 위치 앞당김)는 §3 본문이 p4 좌측 단을 끝까지
% 채워 전폭 하단 공간이 안 나와 실패했다. 넣으려면 표를 \scriptsize 로 줄이거나 §3 본문을
% 줄여야 한다.
% 자동 생성 — tools/make_paper_tables.py. 손으로 고치지 말 것.
\begin{table}[t]
\centering
\caption{교전 시나리오와 관측 설정. 값은 학습 체크포인트의 설정에서 직접 읽었다.}
\label{tab:scenario}
\footnotesize
\resizebox{\ifdim\width>\linewidth\linewidth\else\width\fi}{!}{%
\begin{tabular}{lrl}
\toprule
항목 & 값 & 단위 \\
\midrule
적 USV 수 $M$ & 10 & 척 \\
방어정 수 $P$ & 3 & 척 \\
적 속력 & 9.0 & m/s \\
방어정 속력 & 6.0 & m/s \\
교전 해역 한 변 & 12600 & m \\
그물벽 최대 길이 & 450 & m \\
관측 격자 반폭 & 6300 & m \\
관측 격자 해상도 & 50 & 칸 \\
적 무리 최대 수 $C$ & 4 & 개 \\
관측 픽셀 한 변 & 252 & m \\
방어정당 그물 (학습 / 평가) & 1 / 3 & 장 \\
정책 파라미터 수 & 244,389 & 개 \\
\bottomrule
\end{tabular}
}
\\[3pt]
\begin{minipage}{\linewidth}\footnotesize
방어정 속력이 적 속력보다 낮다 --- 추격이 성립하지 않으므로 접근 회랑의 선제 차단이 유일한 대응이다. 그물 보유량은 학습 시 1장, 평가·배포 시 3장이다: 정책은 결정마다 그물벽 하나를 놓으므로 결정당 행동공간이 보유량과 무관하고, 1장으로 학습한 정책이 재학습 없이 3장 설정에서 동작한다.
\end{minipage}
\end{table}


\begin{table*}[!tb]
\centering
\caption{평가 지표 10종. 방향은 개선이 어느 쪽인지를 나타내며, 기호는 본문에 정의하였다.}
\label{tab:metrics}
\footnotesize
\renewcommand{\arraystretch}{1.25}
\begin{tabular}{llcp{0.27\linewidth}p{0.37\linewidth}}
\toprule
관점 & 지표 & 방향 & 정의 & 뜻 \\
\midrule
\multirow{2}{*}{방어 성과}
 & 포획률           & $\uparrow$   & $\capr = N_\mathrm{cap} / (N_\mathrm{cap} + N_\mathrm{brc})$ & 교전이 결판난 적선 중 포획된 비율 \\
 & 돌파 수          & $\downarrow$ & $N_\mathrm{brc}$ & 모선까지 도달한 적선 수 \\
\midrule
\multirow{2}{*}{안전}
 & 충돌률           & $\downarrow$ & $N_\mathrm{col} / P$ & 방어정 1척당 아군 상호·모선 충돌 건수 \\
 & 그물 걸림        & $\downarrow$ & $N_\mathrm{tng}$ & 아군이 설치된 그물에 얽힌 횟수 \\
\midrule
\multirow{2}{*}{자원 효율}
 & 포획당 그물      & $\downarrow$ & $N_\mathrm{net} / N_\mathrm{cap}$ & 적 1척 포획에 소모한 그물 수 \\
 & 포획당 이동거리  & $\downarrow$ & $\sum_{p=1}^{P} D_p \,/\, N_\mathrm{cap}$ & 적 1척 포획당 팀 총 이동거리 \\
\midrule
기동 비용
 & 총 선회량        & $\downarrow$ & $\sum_{p=1}^{P} \sum_{t} \lvert \Delta\psi_{p,t} \rvert$ & 누적 방위 변화량(궤적 부드러움) \\
\midrule
\multirow{3}{*}{교전 품질}
 & 포획 이격거리    & $\uparrow$   & $\frac{1}{N_\mathrm{cap}} \sum_{k} \lVert \mathbf{z}_k - \mathbf{m} \rVert$ & 포획 지점의 모선 이격 --- 공간축 조기 제압 \\
 & 평균 포획 시각   & $\downarrow$ & $\frac{1}{N_\mathrm{cap}} \sum_{k} t_k$ & 포획이 일어난 시각 --- 시간축 조기 제압 \\
 & 전개 시 적 거리  & $\downarrow$ & $\frac{1}{N_\mathrm{net}} \sum_{d} \min_{m \in \mathcal{A}} \lVert \mathbf{a}_d - \mathbf{e}_m \rVert$ & 그물 전개를 시작한 순간 최근접 적까지 거리 \\
\bottomrule
\end{tabular}
\end{table*}

\begin{figure*}[!t]
\centering
\includegraphics[width=\textwidth]{fig2.png}
\caption{세 가지 공격 포메이션 --- 집중(왼쪽), 양동(가운데), 파상(오른쪽)}
\label{fig:formations}
\end{figure*}

교전 해역 중앙에 고가치 모선 한 척이 있고, 가장자리에서 적 USV $M$척이 모선을 향해
돌진한다. 적은 무작위로 흩어져 오지 않고 세 가지 공격 포메이션 --- 집중, 양동, 파상 ---
가운데 하나를 취한다(Fig.~\ref{fig:formations}; 정의는 아래). 방어정 $P$척은
적선보다 느리며, 각각 한정된 수의 그물벽을 전개할 수 있다. 구체적인 값은
Table~\ref{tab:scenario}에 정리하였다.

방어정은 적을 따라가는 대신 적이 지나갈 자리를 예측하여 선점해야 한다. 그물벽은 두
좌표를 잇는 선분으로 정의되며, 방어정이 그 선분을 따라 이동하면서 전개한다. 물리적으로는
방어정 후미의 전개 장치가 부유식 그물을 해면에 풀어 놓고, 그 위를 지나는 적 USV의
프로펠러가 그물에 얽혀 추진을 잃는 것을 포획으로 본다(Fig.~\ref{fig:netconcept}). 전개된 그물은
해면에 고정되지 않아, 시간이 지날수록 차단 효과가 감소하며, 미리 전개된
그물은 적에게 노출되어 회피의 여지를 준다. 따라서 그물벽은 적의 도달 직전, 적에
가까운 위치에 전개하는 것이 유리함을 가정한다.

\begin{figure}[H]
\centering
\includegraphics[width=0.7\columnwidth]{fig3.png}
\caption{후미 투하식 부유 그물에 의한 포획 개념 --- 그물 위를 지나는 적 USV의 프로펠러가
         얽혀 추진을 잃는다.}
\label{fig:netconcept}
\end{figure}

세 포메이션의 정의는 다음과 같다.
\begin{itemize}[leftmargin=*, itemsep=3pt, topsep=3pt]
\item \textbf{집중(concentrated)}: 적 전 병력이 한 방위에 모여 동시에 접근한다. 차단해야
      할 접근로의 방위가 동일하므로, 소수의 방어정으로도 대응이 가능하다.
\item \textbf{양동(diversionary)}: 적이 세 방위로 나뉘어 \emph{동시에} 접근한다. 접근로가
      셋으로 갈라지므로 가용 방어정을 모두 투입해야 하며, 어느 방어정을 어느 무리에 붙이는
      배정과 그 경로의 교차 여부가 결과를 좌우한다.
\item \textbf{파상(wave)}: 적이 3개 단(rank)으로 나뉘어 단마다 방위를 조금씩 달리하며
      거리차를 두고 순차적으로 도달한다. 앞 단을 막은 자원이 뒤 단에도 쓰일 수 있어야 하므로 그물벽을 \emph{어느
      자리에} 놓는지가 다음 단의 차단 가능성을 좌우하며, 앞 단에 소모된 그물벽을 대신할
      능동적 재배치가 방어율을 결정한다.
\end{itemize}

\subsection{평가 지표}
\label{sec:metrics}

방어 성과의 지표는 다섯 관점에서 총 10개 지표를 정의한다(Table~\ref{tab:metrics}). 지표들은 에피소드
단위로 계산되며, 표의 정의에서 $N_\mathrm{cap}$·$N_\mathrm{brc}$ 는 포획된·모선까지 돌파한 적선
수, $N_\mathrm{col}$ 은 아군 상호 및 모선 충돌 건수, $N_\mathrm{tng}$ 는 아군이 설치된
그물에 얽힌 횟수, $N_\mathrm{net}$ 은 소모한 그물 수, $D_p$ 와 $\Delta\psi_{p,t}$ 는
방어정 $p$ 의 총 이동거리와 step $t$ 의 방위 변화량, $\mathbf{m}$ 은 모선 위치,
$\mathbf{z}_k$ 와 $t_k$ 는 $k$ 번째 포획 지점과 시각, $\mathbf{a}_d$ 는 $d$ 번째 그물
전개를 시작한 방어정의 위치, $\mathcal{A}$ 는 그 시점의 생존 적선 집합, $\mathbf{e}_m$ 은 적선 $m$ 의 위치이다.
자원 효율 지표는 교전하지 않은 조건이 가장 효율적으로 평가되는 것을 막기 위해 총량이
아니라 포획당 비율로 둔다. 교전 품질의 지표들은 거리축과 시간축으로 나누어, 모선에서 먼 곳의
포획과 이른 포획을 함께 본다.

% 기호표는 §3 에서 기호가 정의되므로 여기(§3 끝)에 둔다. §4 첫머리에 두면 p6 상단을
% 차지해 Fig.2/3 을 밀어낸다(2026-09-16 배치 실측). 본문 인용은 §7 뿐이라 위치가 자유롭다.
%-------------------------------------------------- 기호표
%  ★ 이 표에 기호를 추가하려면 본문에 실제로 그 기호가 있는지 먼저 확인할 것.
%    표에만 있고 본문에 없는 기호는 독자를 찾아 헤매게 만든다.
\begin{table*}[!tp]
\caption{주요 기호. 혼동되기 쉬운 수량은 서로 다른 문자를 쓴다 --- 특히 무리 수 $C$,
         경유점 수 $K$, 학습의 그룹 후보 수 $G$ 는 구분한다.}
\label{tab:nomenclature}
% 다른 표와 같은 \footnotesize 로 둔다. \small(10.95pt)로 두면 자연 폭이 커져
% \resizebox 가 발동하고, 그러면 글자가 오히려 더 작아진다(측정: 10.95 -> 9.53).
\centering\footnotesize
% 판면보다 넓을 때만 줄인다(2단 양단걸침 ≈484pt vs elsarticle preprint 390pt).
\resizebox{\ifdim\width>\linewidth\linewidth\else\width\fi}{!}{%
\begin{tabular}{@{}ll@{\hspace{1.2em}}ll@{}}
\toprule
\textbf{기호} & \textbf{뜻} & \textbf{기호} & \textbf{뜻} \\
\midrule
$M$            & 적 USV 수                            & $\mathbf{o}_p$ & 방어정 $p$ 의 자체 상태 벡터 \\
$P$            & 방어정 수                            & $\mathbf{q}_p$ & 방어정 $p$ 의 질의 벡터 \\
$C$, $\hat{C}$ & 무리 수 상한, 활성 무리 수            & $\mathbf{f}_i$ & 픽셀 $i$ 의 키 벡터 \\
$K$            & 방어정 1척이 지목하는 경유점 수      & $s_{p,i}$      & 방어정 $p$ 가 픽셀 $i$ 에 준 점수 \\
$G$            & 학습의 그룹 후보 수                  & $(\boldsymbol{\gamma}_{\mathrm{F}},\boldsymbol{\beta}_{\mathrm{F}})$ & FiLM 배율·이동 계수 \\
\addlinespace[2pt]
$\tau_j$       & 무리 $j$ 의 위협도                   & $R_k$          & 후보 $k$ 의 팀 보상 \\
$\vartheta_m$, $g_{(q)}$, $\delta_{\mathrm{gap}}$ & 적선 방위, 정렬 방위 간격, 분할 임계 & $\mathbf{c}_j$, $n_j$ & 무리 $j$ 의 중심, 척수 \\
$d_{pj}$       & 방어정 $p$ $\to$ 무리 $j$ 요격점 거리 & $A_k$         & 후보 $k$ 의 그룹 상대 이득 \\
$\beta$        & 배정 유지 보너스 (식~\ref{eq:sticky}) & $\Phi$        & 위협 근접도 퍼텐셜 \\
$\lambda_{\mathrm{L}}$ & 잉여 방어정의 층 분리 증분           & $\gamma$       & 퍼텐셜 성형 할인 계수 \\
\addlinespace[2pt]
$\mathbf{I}$   & 요격점                               & $\mathcal{V}^{(k)}_p$ & $k$ 번째 선택 시점의 후보 픽셀 집합 \\
$\mathbf{t}_k$ & $k$ 번째 목표점 (연속 좌표)          & $r_{\mathrm{dup}}$ & 경유점 중복 배제 반경 \\
$\hat{\mathbf{r}}$ & 모선 $\to$ 요격점 단위 벡터      & $L_{\mathrm{net}}$ & 그물벽 최대 길이 \\
$\mathbf{a}_p$ & 방어정 $p$ 의 위치                   & $\dz$          & 대응표본 효과크기 \\
\addlinespace[2pt]
$E$, $\ell$    & 관측 반폭, 픽셀 크기                 & $H$           & 후보 평가 창 길이 \\
$\iota(\cdot)$ & 월드 $\to$ 픽셀 인덱스 사상 (식~\ref{eq:world2pix}) & $\eta_t$, $\lambda_{\mathrm{H}}(t)$ & 학습률·엔트로피 계수 스케줄 \\
$\mathbf{u}_k$, $\boldsymbol{\sigma}_o$ & 부화소 오프셋과 그 표준편차 & $w_n$ & 유효 월드 마스크 \\
\bottomrule
\end{tabular}
}
\\[3pt]
\begin{minipage}{\linewidth}\footnotesize
기호의 실험 설정값은 Table~\ref{tab:scenario}와 각 정의 위치의 본문에 있다.
\end{minipage}
\end{table*}
